% Source: docs/blueprint/01-research/01-sota-landscape.md; comparison axes only, no measured numbers.
\section{Background}
\label{sec:background}

Four lineages of recursive self-improvement are relevant. Archive-based self-modification of a coding agent evaluates whole-scaffold rewrites on a benchmark and keeps parents by score. Evolutionary program search applies population operators to code under a fixed evaluator. The autoresearch loops edit a single training file under a fixed wall-clock budget and keep a change if a validation metric improves. Self-play weight trainers optimise adapters against verifiable rewards. Table~\ref{tab:related} contrasts them on the axes that the present design changes: where the posterior over effects lives, whether evidence carries provenance, and whether the evaluator is reachable by the mutated process.

\begin{table}[t]
\centering
\caption{Recursive self-improvement lineages compared on the three components Pravrudhi changes. ``Reachable'' means the process being mutated can, in principle, write to the evaluator's inputs or outputs.}
\label{tab:related}
\begin{tabular}{lccc}
\toprule
\textbf{Lineage} & \textbf{Posterior} & \textbf{Typed evidence} & \textbf{Evaluator reachable} \\
\midrule
Archive self-modification & none (score archive) & no & yes \\
Evolutionary program search & none (population) & no & yes \\
Autoresearch keep-if-better & none & no & yes \\
Self-play adapter training & implicit in reward & no & partly \\
Pravrudhi & explicit, hierarchical, external & yes & no \\
\bottomrule
\end{tabular}
\end{table}
